\section{Bodegas y Centros de Distribucion}

\subsection{Profundidad optima de pallets}

\begin{ejercicio}{Guia 9, problema 1}
Se disena una bodega de pallets a piso. Los pasillos miden $4.2$ m de ancho y los pallets miden $1.2$ m por $1.0$ m. La cara mas angosta queda hacia el pasillo. Los datos son:
\[
\begin{array}{c|ccc}
\text{SKU} & q_i & z_i & T_i\\
\hline
A&32&1&5\\
B&28&2&4\\
C&14&3&6
\end{array}
\]
donde $q_i$ son pallets pedidos, $z_i$ es altura de apilado en pallets y $T_i$ es ciclo de pedido en semanas. Se pide profundidad optima por SKU y una profundidad unica para todos.
\end{ejercicio}

\begin{materia}{profundidad dedicada}
Para cada SKU $i\in I=\{A,B,C\}$:
\[
  d_i^\star=\sqrt{\frac{a q_i}{2 z_i}},
\]
donde:
\begin{itemize}
  \item $d_i^\star$ es la profundidad optima en numero de pallets.
  \item $a$ es la razon entre ancho de pasillo y ancho efectivo del pallet usado en el frente.
  \item $q_i$ es la cantidad de pallets recibida en una reposicion.
  \item $z_i$ es la altura de apilado.
\end{itemize}
Como la guia indica que la cara mas angosta queda hacia el pasillo, usa:
\[
  a=\frac{4.2}{1.2}=3.5.
\]
\end{materia}

\begin{solucion}{profundidad por SKU}
\[
d_A^\star=\sqrt{\frac{3.5(32)}{2(1)}}=7.48,
\qquad
d_B^\star=\sqrt{\frac{3.5(28)}{2(2)}}=4.95,
\qquad
d_C^\star=\sqrt{\frac{3.5(14)}{2(3)}}=2.86.
\]
Las profundidades son numeros de pallets. En la practica se redondea considerando restricciones fisicas y politica operativa; matematicamente estos son los optimos continuos.
\end{solucion}

\begin{materia}{profundidad unica}
Si todos los SKU deben compartir una misma profundidad $d$, se pondera por la frecuencia relativa de reposicion. Cuando los ciclos son distintos, se toma un horizonte comun $H$ igual al minimo comun multiplo de los ciclos:
\[
  H=\operatorname{mcm}(5,4,6)=60.
\]
La cantidad de reposiciones de SKU $i$ dentro del horizonte es:
\[
  r_i=\frac{H}{T_i}.
\]
La guia usa:
\[
  d^\star=
  \sqrt{
  \frac{a}{2H}
  \sum_{i\in I}r_i\frac{q_i}{z_i}
  }.
\]
\end{materia}

\begin{solucion}{profundidad comun}
Las frecuencias relativas son:
\[
  r_A=\frac{60}{5}=12,\qquad
  r_B=\frac{60}{4}=15,\qquad
  r_C=\frac{60}{6}=10.
\]
Entonces:
\[
  d^\star=
  \sqrt{
  \frac{3.5}{2\cdot 60}
  \left(
  12\frac{32}{1}
  +15\frac{28}{2}
  +10\frac{14}{3}
  \right)
  }
  =
  4.3227.
\]
\end{solucion}

\begin{alerta}{supuestos que pueden preguntar}
La formula no considera restricciones fisicas de la bodega, usa un ancho unico de pasillo, trabaja con pallets completos y supone que las tasas relativas de venta se mantienen estables. Si te piden ventajas/desventajas de almacenamiento compartido, la idea central es espacio mejor utilizado a cambio de mayor complejidad operativa y necesidad de sistema de informacion.
\end{alerta}
